\section{Trabalhos Relacionados}
\label{sec:trabalhos_relacionados}

Os fundamentos do treinamento de redes neurais multicamadas através da descida de gradiente e retropropagação foram consolidados no trabalho seminal de Rumelhart et al. \cite{rumelhart1986learning}. Desde então, diversos arcabouços computacionais foram propostos para simplificar o cálculo automático de gradientes e a construção de topologias complexas.

Na literatura recente, destacam-se abordagens voltadas à interpretabilidade e representação explícita de grafos de computação, as quais auxiliam pesquisadores e estudantes a diagnosticar fenômenos como o desvanecimento de gradientes (\textit{vanishing gradients}) e a saturação de funções de ativação não lineares.
